\documentclass[10pt,a4paper]{amsart}
\usepackage[margin=19mm]{geometry}
\usepackage{amsmath,amssymb,mathtools}
\usepackage[scr=rsfs,cal=euler]{mathalfa}
\usepackage{xfrac}
\usepackage[unicode,hidelinks,bookmarks=false]{hyperref}
\newcommand{\ie}{i.e.\ }
\newcommand{\cf}{cf.\ }
\newcommand{\resp}{respectively\ }
\newcommand{\Subsection}{\subsection}
\providecommand{\nobreakdash}{\nobreak\hbox{-}}
\setlength{\emergencystretch}{2em}
\multlinegap0pt
\usepackage{amssymb}
\usepackage[shortlabels]{enumitem}
\usepackage[normalem]{ulem}
\usepackage{xcolor}
\newtheorem{lemma}{Lemma}
\title{A note on quasi-perfect morphisms}
\author{Timothy De Deyn, Pat Lank, Kabeer Manali Rahul}
\date{Source: arXiv:2508.10845v3 (CC BY 4.0)}
\begin{document}
\maketitle

\noindent\emph{Source and licence.} A note on quasi-perfect morphisms. Version arXiv:2508.10845v3 (CC BY 4.0), distributed by arXiv under the Creative Commons Attribution 4.0 International licence, http://creativecommons.org/licenses/by/4.0/, as recorded in the arXiv licence metadata for this exact version and shown on the abstract page https://arxiv.org/abs/2508.10845v3. Full licence notice: \texttt{licenses/arxiv-2508.10845-CC-BY-4.0.txt}.\par\smallskip

\noindent\textit{Editorial scope.} Counted unit: the article's lemma lem:perfect\_morphism\_via\_generation (source0.tex lines 382--385). The article's complete original proof of it is delivered with it (source0.tex lines 387--394, one \texttt{proof} environment of 8 lines). No mathematics is written, repaired or re-derived: the statement and the proof are the article's own lines, in the article's own notation, and no retained line is substituted or re-flowed.  No further source-local context is needed: every cross-reference of every retained line resolves inside the two delivered spans.  Nothing was omitted from any retained span, so the package carries no omission record.  No retained line contains a citation, so the wrapper carries no bibliography and no bibliography entry.  No retained line contains a figure environment, an image-inclusion command or drawing code, so no image is delivered and the package declares no asset dependency.  Editorial formatting only, in addition: an AMS \texttt{amsart} preamble that restates the article's own declarations and macros used by the retained lines, namely the environments \texttt{lemma} on separate counters, together with the standard packages those lines need.  The retained environments print in this document as Lemma 1 in order of appearance, whereas the article prints the same items with the numbering of its own full document.  The complete native source of the article as distributed in the arXiv e-print for this exact version is bundled unchanged as \texttt{sources/183284/source0.tex}.
\begin{lemma}
    \label{lem:perfect_morphism_via_generation}
    A morphism $f\colon Y \to X$ of quasi-compact quasi-separated algebraic spaces is quasi-perfect if and only if $\mathbf{R}f_\ast P\in \operatorname{Perf}(X)$ for any compact generator $P$ of $Y$.
\end{lemma}

\begin{proof}
    If $f$ is quasi-perfect, it is clear that $\mathbf{R}f_\ast P\in \mathbf{R}f_\ast \operatorname{Perf}(Y)\subseteq \operatorname{Perf}(X)$.
    For, the converse direction, let $P$ be compact generator on $Y$. Our hypothesis is $\mathbf{R}f_\ast P\in \operatorname{Perf}(X)$. 
    The desired claim then follows from the following string of inclusions:
    \begin{displaymath}
        \mathbf{R}f_\ast \operatorname{Perf}(Y) = \mathbf{R}f_\ast \langle P \rangle \subseteq \langle \mathbf{R}f_\ast P \rangle \subseteq \operatorname{Perf}(X).\qedhere
    \end{displaymath}
\end{proof}

\end{document}
